\documentclass[10pt,a4paper]{amsart}
\usepackage[margin=19mm]{geometry}
\usepackage{amsmath,amssymb,mathtools}
\usepackage[scr=rsfs,cal=euler]{mathalfa}
\usepackage{xfrac}
\usepackage[unicode,hidelinks,bookmarks=false]{hyperref}
\newcommand{\ie}{i.e.\ }
\newcommand{\cf}{cf.\ }
\newcommand{\resp}{respectively\ }
\newcommand{\Subsection}{\subsection}
\providecommand{\nobreakdash}{\nobreak\hbox{-}}
\setlength{\emergencystretch}{2em}
\multlinegap0pt
\usepackage{mathtools}
\usepackage{amssymb}
\newtheorem{thm}{Theorem}
\DeclareMathOperator{\BunConn}{BunConn}
\newcommand{\C}{\mathbb C}
\newcommand{\E}{\mathcal E}
\DeclareMathOperator{\GL}{GL}
\newcommand{\PP}{\mathbb P}
\DeclareMathOperator{\Res}{Res}
\newcommand{\XX}{\mathbb X}
\newcommand{\bw}{\mathbf w}
\DeclareMathOperator{\dimv}{\underline\dim}
\DeclareMathOperator{\rank}{rank}
\DeclareMathOperator{\rep}{rep}
\title{The Deligne--Simpson Problem}
\author{William Crawley-Boevey, Andrew Hubery}
\date{Source: arXiv:2509.11998v5 (CC BY 4.0)}
\begin{document}
\maketitle

\noindent\emph{Source and licence.} The Deligne--Simpson Problem. Version arXiv:2509.11998v5 (CC BY 4.0), distributed by arXiv under the Creative Commons Attribution 4.0 International licence, http://creativecommons.org/licenses/by/4.0/, as recorded in the arXiv licence metadata for this exact version and shown on the abstract page https://arxiv.org/abs/2509.11998v5. Full licence notice: \texttt{licenses/arxiv-2509.11998-CC-BY-4.0.txt}.\par\smallskip

\noindent\textit{Editorial scope.} Counted unit: the article's thm thm (source0.tex lines 857--865). The article's complete original proof of it is delivered with it (source0.tex lines 867--889, one \texttt{proof} environment of 23 lines). No mathematics is written, repaired or re-derived: the statement and the proof are the article's own lines, in the article's own notation, and no retained line is substituted or re-flowed.  No further source-local context is needed: every cross-reference of every retained line resolves inside the two delivered spans.  Nothing was omitted from any retained span, so the package carries no omission record.  No retained line contains a citation, so the wrapper carries no bibliography and no bibliography entry.  No retained line contains a figure environment, an image-inclusion command or drawing code, so no image is delivered and the package declares no asset dependency.  Editorial formatting only, in addition: an AMS \texttt{amsart} preamble that restates the article's own declarations and macros used by the retained lines, namely the environments \texttt{thm} on separate counters, together with the standard packages those lines need.  The retained environments print in this document as Theorem 1 in order of appearance, whereas the article prints the same items with the numbering of its own full document.  The complete native source of the article as distributed in the arXiv e-print for this exact version is bundled unchanged as \texttt{sources/185074/source0.tex}.
\begin{thm}
We fix the data $\alpha\in\Gamma_\bw$, $\xi=(\xi_{ip})$ and $\zeta=(\zeta_{ip})$ as above, yielding tuples of conjugacy classes $(C'_1,\ldots,C'_k)$ and $(C_1,\ldots,C_k)$ respectively.

(i) There exists a representation $\pi_1\to\GL_n(\C)$ with the image of $g_i$ lying in $\overline C'_i$ for all $i$ if and only if there exists $(\E,\nabla)\in\BunConn_\zeta\XX$ with $\dimv\E=\alpha$.

(ii) There exists a representation $\pi_1\to\GL_n(\C)$ with the image of $g_i$ lying in $C'_i$ for all $i$ if and only if there exists a strict $(\E,\nabla)\in\BunConn_\zeta\XX$ with $\dimv\E=\alpha$.

(iii) There exists an irreducible representation $\pi_1\to\GL_n(\C)$ with the image of $g_i$ lying in $C'_i$ for all $i$ if and only if there exists an irreducible $(\E,\nabla)\in\BunConn_\zeta\XX$.
\end{thm}

\begin{proof}
Write $\alpha=n\alpha_\ast+\sum_{i,p}n_{ip}\alpha_{ip}$. We know that $A_i\in\GL_n(\C)$ lies in $\mathcal C'_i$ if and only if
\[ \rank(A_i-\xi_{ip})\cdots(A_i-\xi_{i1}) = n_{ip}. \]
This is equivalent to the existence of a flag
\[ \C^n = V_{i0} \supseteq V_{i1} \supseteq \cdots \supseteq V_{iw_i} = 0, \quad \dim V_{ip}=n_{ip}, \]
with $(A_i-\xi_{ip})(V_{ip-1})=V_{ip}$ for all $p$.

We can similarly characterise the closure $\overline C'_i$ by saying there exists such a flag with $(A_i-\xi_{ip})(V_{ip-1})\subseteq V_{ip}$ for all $p$.

(i) By the Riemann--Hilbert Correspondence, the existence of a representation with $A_i\in\overline C'_i$ for all $i$ is equivalent to the existence of some $(E,\nabla)\in\BunConn_\xi\PP^1$ with $\Res_{a_i}\nabla\in\overline C_i$ for all $i$. By the above observation, this is equivalent to the existence of a certain flag $(E_{ip})$ in each of the fibres $E_{a_i}$, in which case $\E=(E,E_{ip})$ is a parabolic bundle on $\XX$ having dimension vector $\alpha$, and $\nabla$ is a $\zeta$-connection on $\E$.

(ii) Again, having a representation with $A_i\in C'_i$ for all $i$ is equivalent to finding some $(E,\nabla)\in\BunConn_\xi\PP^1$ with $\Res_{a_i}\nabla\in C_i$ for all $i$, which in turn is equivalent to the existence of a certain flag $(E_{ip})$ in each of the fibres $E_{a_i}$. Then $\E=(E,E_{ip})$ is a parabolic bundle on $\XX$ having dimension vector $\alpha$, and $\nabla$ is a $\zeta$-connection on $\E$ such that the pair $(\E,\nabla)$ is strict.

We observe that $(\E,\nabla)$ is precisely the image of $(E,\nabla)$ under the left adjoint.

(iii) Under the Riemann--Hilbert Correspondence, irreducible objects in $\rep_\xi\pi_1$ correspond to irreducible objects in $\BunConn_\zeta\PP^1$. The claim therefore follows provided both the forgetful functor $F$ and its left adjoint $L$ send irreducibles to irreducibles.

Now, it is clear from our explicit description of $L$ that the unit $\mathrm{id}\to FL$ is an isomorphism, and hence that $L$ is fully faithful. Also, $L$ preserves monomorphisms and the counit $LF\to\mathrm{id}$ is a monomorphism. Finally, as a right adjoint, $F$ necessarily preserves monomorphisms.

Thus, if $(E,\nabla)$ is irreducible, and $(\E,\nabla)=L(E,\nabla)$, then any non-zero subobject $(\E',\nabla')$ of $(\E,\nabla)$ maps to a non-zero subobject of $(E,\nabla)$, so equals $(E,\nabla)$. The counit then yields an inclusion $(\E,\nabla)\subseteq(\E',\nabla')$, which must therefore be an isomorphism, and hence $(\E,\nabla)=L(E,\nabla)$ is irreducible in $\BunConn_\zeta\XX$.

Conversely, if $(\E,\nabla)$ is irreducible, then it must be equal to its subobject $LF(\E,\nabla)$, so $(\E,\nabla)$ is strict. If $(E',\nabla')$ is a non-zero subobject of $(E,\nabla)\coloneqq F(\E,\nabla)$, then it is sent under $L$ to a subobject $L(E',\nabla')\subseteq(\E,\nabla)$, as $L$ preserves monomorphisms, and hence we get equality. Using that the unit is an isomorphism we see that $(E',\nabla')=(E,\nabla)$, and hence that $(E,\nabla)$ is irreducible.
\end{proof}

\end{document}
